%% Chapter 8 - Commercialisation and Sustainability
%% Sections: 8.1 Value proposition and target users; 8.2 Market and competitive
%% position; 8.3 Business or sustainability model; 8.4 Scalability, intellectual
%% property and compliance; 8.5 Contribution to the Sustainable Development Goals.
%% Master rules: no fabricated market statistics; no claim of legal/certification
%% compliance without evidence; speculative commercialisation written as potential
%% pathways only.

\chapter{Commercialisation and Sustainability}\label{ch:commercialisation}

This chapter examines how the Intent Recognition Engine could be taken from a
research artefact to a deployed capability. It states the value proposition and
the users it would serve, positions the engine against the surveyed landscape,
and sets out plausible---not claimed---commercialisation and sustainability
pathways. Consistent with the evidence-first rule of this dissertation, no
market statistic is fabricated here, no revenue or adoption figure is asserted,
and every legal or certification statement is either grounded in the verified
record of the preceding chapters or explicitly marked as a requirement rather
than an achievement.

\section{Value Proposition and Target Users}\label{sec:value-proposition}

The value proposition follows directly from the measured result of
Chapter~\ref{ch:results}. A security organisation that wants to let analysts
drive tooling through natural language faces a choice between correctness and
automation: a raw language model offers fluency but, as Experiment~1 shows,
transfers every malformed, misclassified or hallucinated output into the actions
a pipeline would execute---a 0\% catch rate and a 32\%/31\% false-positive
rate before guarding. The IRE converts that choice into a combination: on the
evaluated 1,500-record population, the guarded pipeline raises flaw catch from
0\% to 64--72\%, lifts decision accuracy from 69--78\% to 82--86\%, and halves
the false-positive rate (14\%/15\%), while remaining offline, CPU-only,
single-inference and retrieval-free
(Section~\ref{sec:exp1}). The value is therefore not ``an LLM that understands
security prose''---it is ``a head-of-pipeline component that guarantees
structurally valid, schema-constrained, authorised and in-scope intent
contracts despite the underlying model''.

This value concentrates where three conditions coincide: the organisation uses
offensive-security tooling, operates under data- or policy-constraints that
disallow sending prompts to a hosted model or a retrieval service, and has an
automation chain that can consume validated JSON contracts rather than prose
(Section~\ref{sec:component-architecture}). Three user groups follow.

\emph{Security analysts and operators} are the primary users. For them the
engine removes the manual translation stage between intent and tool invocation,
and the safety invariant (Equation~\ref{eq:safety-invariant}) means a
misremembered tool syntax, an out-of-scope target or a role violation is caught
by deterministic predicates rather than by the analyst's vigilance. The
audit-friendly structured response and per-request metrics
(Section~\ref{sec:deployment-operation}) also suit post-action review.

\emph{Security-tooling integrators} are the secondary users. Because the IRE
emits a closed, schema-enforced contract with a documented tool hint and
parameter envelope (nine intent classes, validated target types and explicit
CVE/port/modifier fields), it is an integration point for orchestrators and
planners that would otherwise re-parse prose (Section~\ref{sec:component-architecture}).
Within the broader NeuroShell chain (project reference R26-CS-018,
Section~\ref{sec:overall-project-context}) the IRE is Component~01 and is
designed to be the validated eye of the chain.

\emph{Offline-constrained environments} constitute the differentiating third
user group: air-gapped networks, regulated industries, embedded or edge
deployments and research ranges where external inference is unavailable or
prohibited. The engine's only runtime dependencies are model weights, policies
and an operator API key, and nothing leaves the device at inference time
(Section~\ref{sec:deployment-operation}). This is the segment the comparable
copilots in the survey do not serve (Section~\ref{sec:existing-systems}).

\section{Market and Competitive Position}\label{sec:market-position}

The competitive position is stated relative to the verified
existing-systems survey (Section~\ref{sec:existing-systems}) rather than to
unverifiable market sizing. The survey found a four-way split: deterministic
tooling (for example Nmap and the exploitation framework) that is ground-truth
accurate but requires a human to map intent, generation-side copilots that map
intent but consume human review, integration techniques that connect models to
tools without an intermediate validation stage, and post-hoc guardrail
techniques that check model behaviour after the fact. In those four categories,
any candidate is strong in one and absent in another: none combined
machine-consumed output, schema and scope/authorisation enforcement, and
offline, CPU-only, retrieval-free operation over a closed intent taxonomy
(Section~\ref{sec:gap-synthesis} and Table~\ref{tab:existing-systems}).

The IRE occupies that empty combination, and its defensible position rests on
three properties the evaluation supports rather than on marketing claims.
First, the catch-before-dispatch architecture is a behavioural claim with a
measurement: the ablation record shows the validation stack is the dominant
safety contributor (Gemma $-30$, Qwen $-24$ catch points when removed;
Section~\ref{sec:exp3}). Second, the result is reproducible and
backend-independent on the safety dimension---the identical stack orders the
effects of the same way on two distinct offline models
(Section~\ref{sec:exp2})---which weakens a ``vendor-locked model'' objection.
Third, the offline posture is structural, not incidental: it is the deployment
constraint the design is built around, not a configuration.

Two competitive risks are stated openly. The large-model copilots are strong on
coverage and conversation, so the IRE competes only in the automated,
machine-consumed corner of the space, not against interactive assistants. And
because no named commercial product was included in the survey for lack of
verifiable primary sources (Section~\ref{sec:existing-systems}),
the position here is a review-based gap claim, to be re-verified against named
products at any commercialisation step before a broader claim is made.

\section{Business and Sustainability Model}\label{sec:business-model}

Consistent with the master rules, the following is a set of potential pathways,
not a claimed business. All are sketched as options that a founder or
institution could evaluate, and none asserts a market or revenue number.

\subsection*{Open-core with managed hardening services}

The engine, its data set and its evaluation harness are research artefacts
versioned with the project source
(Section~\ref{sec:reproducibility}), and the model backends are open-weight
and locally runnable. A natural model is open-core: the validated core is
published to establish trust and to grow an integrator base, while the
commercial layer concentrates on what a single-team paper cannot ship---the
hardening and certification work itemised in Section~\ref{sec:ip-compliance}
(bounded latency measurement, closed production scope path, credential
management), integration packs for the principal scanning and exploitation
tools, and a maintained contract-schema migration path.

\subsection*{Direct licensing and private deployment}

An alternative or complement is direct licensing of the component to security
product vendors and internal security platforms, under which the licensee embeds
the IRE as the NLU gateway between their users and their tooling. The
deliverable in this pathway is a library-plus-service with a documented contract
schema and the test suite (506 tests, 76\% statement coverage;
Section~\ref{sec:product-results}) as the acceptance baseline the licensee can
run independently.

\subsection*{Sustainability rather than revenue}

A third reading of ``sustainability'' is operational, not financial. Because the
system is CPU-only, runs on a single device without a hosted inference bill,
and holds no external data dependency at inference time
(Section~\ref{sec:deployment-operation}), it has a low and predictable running
cost. That property is itself the sustainable-deployment argument: the engine
can keep providing a security-gate capability in environments or budgets where
hosted copilot licensing is not viable. Any real costing must be produced from
measured latency and resource profiles, neither of which is asserted here
(Section~\ref{sec:discussion-limitations}).

\section{Scalability, Intellectual Property and Compliance}\label{sec:ip-compliance}

\subsection*{Scalability}

The scalability story is deliberately honest about the current measurement.
Each request is a single offline inference, so throughput on one device is
bounded by the model's inference latency and the semantic-cache hit rate; no
latency figures are quoted because they were not measured in this dissertation
(Section~\ref{sec:discussion-limitations}). What scales by design is the
deployment shape, not the single-node speed: because the component is a
stateless, rate-limited API service holding no per-request state beyond the
session context, more capacity is acquired by adding independent instances
behind a load balancer, each with its own model weights, rather than by
contending on one node (Sections~\ref{sec:component-architecture} and
\ref{sec:deployment-operation}). The semantic cache, the bounded-sampling
agreement gate and the deterministic escalation path
(Section~\ref{sec:related-research}) are the design choices that keep the
per-request cost bounded as usage grows.

\subsection*{Intellectual property}

The intellectual-property position rests on what was authored here and what is
third-party. Authored: the seven-stage pipeline composition, the closed
nine-intent taxonomy, the 22-entry alias-to-CVE resolver, the evaluation data
set and harness, and the ablation methodology used to quantify layer
contribution (Section~\ref{sec:component-architecture} and
Chapter~\ref{ch:results}). These are the defensible original elements, and any
commercialisation should identify them as the owned asset
(Section~\ref{sec:discussion-product}). Third-party: the model backends
(\texttt{gemma4:e4b} and \texttt{qwen2.5-coder:7b}) are open-weight, the OWASP
Top~10 for LLM applications \parencite{owasp2025}, RFC~1918 \parencite{rfc1918}
and the vulnerability-identifier registries \parencite{nistnvd} are external
standards referenced by the design, and the scanning-tool syntaxes the engine
emits follow vendor documentation \parencite{lyon2009nmap,metasploit}
(Section~\ref{sec:standards}). A licensing review is therefore a requirement of
any commercial step, not a completed fact: the statement here is that the
original layers are separable from the third-party dependencies, not that any
licence has been executed.

\subsection*{Compliance}

The verified compliance position is narrow and comfortable. The product
processes no personal data at any point: requests, session context and results
remain on the device and involve no identifiable-individual information
(Section~\ref{sec:ethics}), so it falls outside personal-data protection regimes
as a data-processing matter (Section~\ref{sec:standards}). The research itself
required no institutional research-ethics approval, involved no human
participants and contacted no live system, and holds only model weights and
operator-configured policies as persistent assets
(Section~\ref{sec:ethics}).

Two verified findings bound commercial deployment, and both are stated as
requirements rather than resolved items. First, the shipped configuration uses a
default development API key that bypasses transport authentication when left
untouched, and the ``\texttt{production}'' scope-mode value admitted by the
settings schema falls through to the \texttt{WARN} enforcement path rather than
\texttt{BLOCK}---both must be closed before any deployment outside a controlled
lab (Section~\ref{sec:discussion-product}). Second, institutional confirmation
of the applicable declarations (integrity and AI-tools disclosures) is still
pending at the final AUDIT step (Section~\ref{sec:ethics}). No certification
(metric, standard or regulatory) is claimed as obtained; certification is
positioned as a future hardening task.

\section{Contribution to the Sustainable Development Goals}\label{sec:sdg}

This section states the engine's relationship to the UN Sustainable Development
Goals in the only form the evidence supports: a reasoned contribution
narrative, with no quantified or statistical claim against the SDG indicators.

\emph{SDG~9 (Industry, Innovation and Infrastructure)} is the most direct
mapping. The engine is an applied-AI component for a critical-industry
function, and its deployment argument---validated, offline, low-infrastructure
automation of routine security workflows---speaks to building resilient
cybersecurity capability without depending on large hosted-model infrastructure.
The measured reduction in free-form model error before machine dispatch
(Section~\ref{sec:exp1}) is precisely the reliability improvement that makes an
innovation acceptable in infrastructure automation.

\emph{SDG~16 (Peace, Justice and Strong Institutions)} is served through cyber
resilience. By placing a deterministic, authorisation- and scope-enforced gate
in front of security action, the engine contributes to protecting digital
infrastructure and the institutions that rely on it. The dual-use treatment is
part of that narrative: the system is deployment-gated and scope-enforced, its
failure modes are published rather than suppressed, and its offensive
capability is bounded by the same invariants it enforces
(Section~\ref{sec:ethics}), so the contribution to secure institutions is
conditional on that gate.

\emph{SDG~8 (Decent Work and Economic Growth)} is a secondary mapping: the
system automates the error-prone manual translation of security intent into
tool actions, the kind of task delegation consistent with freeing analysts for
higher-value work. The claim is qualitative only---the engine changes the
distribution of human effort, and its productivity effect is a matter for
measured deployment, not an assertion made here.

No claim is made against SDG~4, 13 or others without a concrete connection, and
no SDG indicator is cited with a number. This section is deliberately more
modest than a typical sustainability claim because the master rule of this
dissertation is that contributions are either measured or explicitly unmeasured.